Neuroscience has established that sleep is not a passive state of shutdown but an "optimal neurophysiological state" for memory consolidation \citep{wilson1994reactivation}. It functions as an "offline" processing phase, allowing neural networks to be "repeatedly reactivated in the absence of new sensory input".

This process is particularly evident in NREM (non-REM) sleep. Groundbreaking rodent studies have shown that patterns of hippocampal place cell activity generated while exploring a maze during wakefulness are "replayed" during subsequent NREM sleep \citep{wilson1994reactivation, pavlides1989placecell}. In humans, brain regions engaged during a learning task are "preferentially activated during sleep," and the strength of this activation predicts subsequent memory improvement \citep{stickgold2000tetris}. This "replay" is not a perfect copy; it occurs in "intermittent bursts" on a faster timescale and, crucially, integrates \textit{multiple related memories}. This integration is thought to facilitate insight, allowing the brain to detect "shortcuts" or novel relationships between data points.